\documentclass[a4paper,11pt]{article}
\usepackage{exam-style-341}
\examinehead{341 农业综合知识三 · 网络技术分章练习}{第01章\quad 计算机网络基本概念与性能指标}
\begin{document}

\unitcover{341 农业综合知识三 · 网络技术分章练习}{第01章\quad 计算机网络基本概念与性能指标}{练习（共 100 分）}{约 60 分钟}{计算机网络的定义与功能；互联网基础结构发展的三个阶段；internet 与 Internet 的区别；计算机网络的分类（按范围、拓扑、交换方式）；电路交换、报文交换与分组交换；性能指标（速率、带宽、吞吐量、时延、时延带宽积、RTT、利用率）；协议与体系结构初步。}

\begin{knowmap}
\textbf{定义与功能：}计算机网络是把地理位置不同、具有独立功能的多台计算机通过通信线路互连，实现资源共享和信息传递的系统；向用户提供的最重要功能是\textbf{连通性}和\textbf{共享}\\[2pt]
\textbf{发展脉络：}教材按互联网基础结构分三阶段——从单个 ARPANET 向互连网发展（1983 年 TCP/IP 成为标准）$\rightarrow$ 三级结构互联网（主干网·地区网·校园网）$\rightarrow$ 多层次 ISP 结构互联网\\[2pt]
\textbf{分类：}按范围 广域网 WAN · 城域网 MAN · 局域网 LAN · 个人区域网 PAN；按拓扑 星形·总线·环形·树形·网状；按交换方式 电路交换·报文交换·分组交换\\[2pt]
\textbf{三种交换：}电路交换（建立连接—通信—释放连接，独占资源，时延小但线路利用率低）；报文交换（以整个报文为单位存储转发，时延大）；分组交换（以分组为单位存储转发、统计复用，高效灵活迅速可靠，但有首部开销和排队时延）\\[2pt]
\textbf{性能指标（7 个）：}速率 · 带宽 · 吞吐量 · 时延（发送 + 传播 + 处理 + 排队）· 时延带宽积 · 往返时间 RTT · 利用率\\[2pt]
\textbf{核心公式：}发送时延 $=$ 数据长度 $/$ 发送速率；传播时延 $=$ 信道长度 $/$ 传播速率；时延带宽积 $=$ 传播时延 $\times$ 带宽；信道利用率 $=$ 有数据通过的时间 $/$ 总时间
\end{knowmap}

\begin{keyterms}
\textbf{计算机网络}：把地理位置不同、具有独立功能的多台计算机通过通信线路互连，实现资源共享和信息传递的系统。\\
\textbf{互连网（internet）}：通用名词，泛指由多个计算机网络互连而成的网络，协议可任意选择。\\
\textbf{互联网（Internet）}：专有名词，指采用 TCP/IP 协议族、全球最大的特定互连网，前身是 ARPANET。\\
\textbf{速率}：数据的传送速率，也称数据率或比特率，单位 bit/s。\\
\textbf{带宽}：频域上指信号的频带宽度（Hz），时域上指信道能通过的“最高数据率”（bit/s）。\\
\textbf{吞吐量}：单位时间内通过某个网络或信道的实际数据量，受带宽和网络负载限制。\\
\textbf{发送时延}：主机或路由器把数据帧的全部比特发送到链路上所需的时间，等于数据长度除以发送速率。\\
\textbf{传播时延}：电磁波在信道中传播一定距离所需的时间，等于信道长度除以传播速率。\\
\textbf{时延带宽积}：传播时延与带宽的乘积，表示链路上可容纳的比特数，又称以比特为单位的链路长度。\\
\textbf{往返时间 RTT}：从发送方发送数据开始，到发送方收到接收方确认所经历的时间，包括传播、处理、排队和转发时延。\\
\textbf{利用率}：信道利用率指信道有数据通过的时间占总时间的比例；网络利用率是全网络信道利用率的加权平均值。\\
\textbf{电路交换}：通信前先建立端到端的连接，通信中独占该连接资源，通信结束后释放连接。\\
\textbf{报文交换}：以整个报文为传输单位，采用存储转发方式，时延较大。\\
\textbf{分组交换}：把报文划分为分组，以分组为单位存储转发，每个分组独立选择路由。
\end{keyterms}

\begin{contrast}{电路交换 / 报文交换 / 分组交换（每行依次为：核心特征·资源占用·优缺点）}
\citem{电路交换}{先建立连接，再通信，最后释放连接}{通信全程独占端到端信道}{时延小、实时性好；线路利用率低，不适合突发性计算机数据}
\citem{报文交换}{以整个报文为单位存储转发}{逐段占用链路，无需建立连接}{无需预约带宽；时延大（分钟到小时级），对节点存储要求高}
\citem{分组交换}{把报文分组，逐段存储转发，分组独立选路}{逐段、断续占用链路资源}{高效、灵活、迅速、可靠；有首部开销，可能排队和失序}
\end{contrast}

\begin{contrast}{带宽 / 吞吐量（每行依次为：含义·单位·相互关系）}
\citem{带宽}{频域指频带宽度；时域指信道最高数据率}{Hz 或 bit/s}{是理论上限、额定值}
\citem{吞吐量}{单位时间内实际通过的数据量}{bit/s}{受带宽、负载、拥塞等影响，通常小于带宽}
\end{contrast}

\begin{contrast}{发送时延 / 传播时延（每行依次为：产生位置·计算式·影响因素）}
\citem{发送时延}{机器内部的发送器（网络适配器）}{数据长度 $/$ 发送速率}{与帧长成正比、与发送速率成反比，与链路长度无关}
\citem{传播时延}{机器外部的传输信道媒体}{信道长度 $/$ 传播速率}{与距离成正比、与发送速率无关，取决于传输媒体}
\end{contrast}

\qtypename{一、单项选择题}{10 题，每题 3 分，共 30 分}
\q 计算机网络向用户提供的最重要的功能是（\quad）
\choices{A. 连通性和共享}{B. 提高计算机的运算速度}
{C. 扩大计算机的存储容量}{D. 降低通信费用}

\q 互联网基础结构发展的第一阶段是（\quad）
\choices{A. 从单个网络 ARPANET 向互连网发展}{B. 建成三级结构的互联网}
{C. 形成多层次 ISP 结构的互联网}{D. 进入移动互联网时代}

\q 关于 internet 与 Internet 的说法，正确的是（\quad）
\choices{A. 小写 internet 泛指多个计算机网络互连而成的网络}{B. 大写 Internet 泛指任意互连网}
{C. 二者含义完全相同，可以混用}{D. 小写 internet 必须使用 TCP/IP 协议}

\q 按覆盖范围由大到小排列，正确的是（\quad）
\choicesline{A. 广域网、城域网、局域网、个人区域网}{B. 局域网、城域网、广域网、个人区域网}
{C. 城域网、广域网、局域网、个人区域网}{D. 个人区域网、局域网、城域网、广域网}

\q 电路交换通信的三个阶段依次是（\quad）
\choices{A. 建立连接、通信、释放连接}{B. 存储、转发、确认}
{C. 采样、量化、编码}{D. 封装、传输、解封装}

\q 下列不属于分组交换主要优点的是（\quad）
\choicesline{A. 高效}{B. 灵活}{C. 迅速}{D. 通信中独占端到端信道}

\q 发送时延的计算公式是（\quad）
\choices{A. 数据长度 $/$ 发送速率}{B. 信道长度 $/$ 传播速率}
{C. 数据长度 $\times$ 带宽}{D. 传播时延 $\times$ 带宽}

\q 传播时延与下列哪个因素无关（\quad）
\choicesline{A. 信道长度}{B. 传播速率}{C. 发送速率}{D. 传输媒体}

\q 时延带宽积表示（\quad）
\choices{A. 链路上可容纳的比特数}{B. 网络每秒处理的分组数}
{C. 链路的物理长度（米）}{D. 网络的误码率}

\q 关于往返时间 RTT，下列说法正确的是（\quad）
\choices{A. 只包括发送时延}{B. 包括链路传播时间、末端处理时间、路由器排队和处理时间}
{C. 与网络拥塞程度无关}{D. 等于单向传播时延}

\qtypename{二、多项选择题}{3 题，每题 4 分，共 12 分；每题有两个或两个以上正确项}
\q 计算机网络的性能指标包括（\quad）
\choicesline{A. 速率}{B. 带宽}{C. 吞吐量}{D. 显示器分辨率}

\q 关于三种交换方式，下列说法正确的有（\quad）
\choices{A. 电路交换通信前必须建立连接}{B. 报文交换以整个报文为传输单位}
{C. 分组交换的每个分组可以独立选择路由}{D. 电路交换的线路利用率高，适合传送突发性计算机数据}

\q 关于带宽与吞吐量，下列说法正确的有（\quad）
\choices{A. 带宽可以指信号的频带宽度，单位是 Hz}{B. 计算机网络中的带宽常指最高数据率，单位是 bit/s}
{C. 吞吐量是单位时间内实际通过的数据量}{D. 吞吐量在任何情况下都等于带宽}

\qtypename{三、判断题}{6 题，每题 2 分，共 12 分；正确的画“√”，错误的画“×”并写出改正后的表述}
\q 计算机网络的“资源”只包括硬件，不包括软件和数据。\quad（\quad）\quad 改正：\fixblank

\q 小写 internet（互连网）是一个通用名词，泛指由多个计算机网络互连而成的网络。\quad（\quad）\quad 改正：\fixblank

\q 电路交换在通信过程中独占端到端通信资源，但用它传送突发性的计算机数据时线路利用率往往很低。\quad（\quad）\quad 改正：\fixblank

\q 分组交换采用存储转发方式，分组在各路由器需要排队，可能带来排队时延。\quad（\quad）\quad 改正：\fixblank

\q 发送时延与信道长度有关，传播时延与发送速率有关。\quad（\quad）\quad 改正：\fixblank

\q 信道利用率越高越好，应尽量接近 100\%。\quad（\quad）\quad 改正：\fixblank

\qtypename{四、计算题}{1 题，共 6 分}
\q 主机 A 通过一条带宽为 10 Mb/s、长度为 2000 km 的链路向主机 B 发送一个 5000 bit 的分组，信号在链路中的传播速率为 $2\times10^{8}$ m/s，忽略处理时延和排队时延，确认帧长度忽略不计。求：（1）发送时延；（2）传播时延；（3）往返时间 RTT；（4）时延带宽积。

\qtypename{五、简答题}{2 题，每题 10 分，共 20 分}
\q 简述计算机网络的定义与两大功能，并按交换方式说明电路交换、报文交换、分组交换的主要区别。\anslines{2.4cm}

\q 简述计算机网络的 7 个性能指标，并写出时延的四个组成部分及其各自的影响因素。\anslines{2.4cm}

\newpage
\qtypename{六、综合题}{1 题，共 20 分}
\q 某农业院校的校园网通过出口路由器连接互联网。校园网内部链路速率均为 100 Mb/s，出口路由器到互联网服务提供商的链路长 1000 km、带宽 100 Mb/s，信号传播速率为 $2\times10^{8}$ m/s。现要把一个 200 MB（按 $1\text{ MB}=10^{6}$ B 计）的教学视频从校内服务器上传到云端，忽略校园网内部链路时延与处理时延。请回答：（1）计算该视频的比特数以及出口链路上的发送时延；（2）若管理员认为“链路带宽 100 Mb/s，所以 200 MB 视频只需 16 s 就能传完”，这一说法是否严谨？为什么？（3）计算出口链路的传播时延和时延带宽积，并说明时延带宽积的物理意义；（4）从提高文件传输吞吐量的角度提出两条改进建议。\anslines{5cm}

\end{document}
